\documentclass[11pt,a4paper]{article}

\usepackage[margin=1in]{geometry}
\usepackage{amsmath,amssymb,amsthm}
\usepackage{mathtools}
\usepackage{booktabs}
\usepackage{enumitem}
\usepackage[T1]{fontenc}
\usepackage{lmodern}
\usepackage[protrusion=true,expansion=false]{microtype}
\usepackage[hidelinks]{hyperref}
\usepackage{xcolor}

\newcommand{\R}{\mathbb{R}}
\newcommand{\E}{\mathbb{E}}
\newcommand{\Prob}{\mathbb{P}}
\newcommand{\Ind}[1]{\mathbf{1}\!\left\{#1\right\}}
\newcommand{\VaR}{\mathrm{VaR}}
\newcommand{\ES}{\mathrm{ES}}
\newcommand{\RV}{\mathrm{RV}}
\newcommand{\BV}{\mathrm{BV}}
\newcommand{\RQ}{\mathrm{RQ}}
\newcommand{\TQ}{\mathrm{TQ}}
\newcommand{\Fcal}{\mathcal{F}}
\newcommand{\Mcal}{\mathcal{M}}
\newcommand{\eqdef}{\vcentcolon=}

\newcommand{\intuition}[1]{\medskip\noindent\textit{Intuition.}\ #1\medskip}

\title{\textbf{cryptorisk: mathematical methodology}\\[2pt]
\large Volatility and tail-risk models for crypto assets, and how they are compared}
\author{v2 project documentation}
\date{\today}

\begin{document}
\maketitle

\begin{abstract}
This note collects the mathematics implemented in the \texttt{cryptorisk} code
base: the notation and the out-of-sample protocol
(\S\ref{sec:setup}); the realized measures computed from 5-minute data
(\S\ref{sec:realized}); the sixteen forecasting models, grouped by family, each
with its defining equations and a one-paragraph reading of what it is doing
(\S\ref{sec:models}); and the evaluation battery --- coverage tests, the
Dynamic Quantile test, Expected-Shortfall backtests, strictly consistent
scoring functions, the Diebold--Mariano test and the Model Confidence Set, and
the Berkowitz density test (\S\ref{sec:eval}). It is meant to be read alongside
\texttt{docs/V2\_PLAN.md} (scope) and the model docstrings. Emphasis is on being
explicit and intuitive rather than on proofs, which are in the cited papers.
\end{abstract}

\tableofcontents

%======================================================================
\section{Setup and notation}\label{sec:setup}

\paragraph{Returns.}
Let $P_t$ be the daily closing price of an asset and
$r_t = \ln P_t - \ln P_{t-1}$ its log-return. We work under the physical
(real-world) measure: the goal is to describe what actually tends to happen,
not to price derivatives.

\paragraph{Risk measures.}
For a tail probability $\alpha \in (0,\tfrac12)$ (we use $\alpha \in
\{0.025,\,0.01\}$; FRTB uses $97.5\%$ ES), the one-day-ahead
\emph{Value-at-Risk} and \emph{Expected Shortfall} are
\begin{equation}
  \VaR_{\alpha,t} \eqdef \inf\{x : \Prob(r_{t+1} \le x \mid \Fcal_t) \ge \alpha\},
  \qquad
  \ES_{\alpha,t} \eqdef \E\!\left[\, r_{t+1} \,\middle|\, r_{t+1} \le \VaR_{\alpha,t},\ \Fcal_t \right].
\end{equation}
Both are \emph{return levels}, normally negative. A \emph{violation} (or
exceedance) on day $t+1$ is the event $\{r_{t+1} < \VaR_{\alpha,t}\}$, with
indicator $I_{t+1} \eqdef \Ind{r_{t+1} < \VaR_{\alpha,t}}$. Under a correct
model $\E[I_{t+1}\mid\Fcal_t] = \alpha$ and $\{I_t\}$ is i.i.d.\ Bernoulli$(\alpha)$.

\paragraph{The predictive distribution.}
Every model outputs, each day, a one-step predictive law $\widehat F_{t}$ for
$r_{t+1}$; VaR and ES are its lower-tail quantile and tail mean. Three concrete
shapes cover all sixteen models:
\begin{itemize}[nosep]
  \item \emph{location--scale}: $r_{t+1} = \mu_{t+1} + \sigma_{t+1} Z$ with $Z$
        from a fixed standardized law (unit variance where it exists);
  \item \emph{sample-based}: quantiles of a (possibly weighted) empirical sample,
        optionally rescaled by $\sigma_{t+1}$;
  \item \emph{quantile-only}: a $(\VaR,\ES)$ pair produced directly, with no
        full density (CAViaR, the MS-GARCH bridge).
\end{itemize}
The probability integral transform (PIT) $u_{t+1} = \widehat F_t(r_{t+1})$ is
recorded whenever a full density is available.

\paragraph{Walk-forward protocol.}
Let the sample run $t = 1,\dots,T$ and fix an out-of-sample start $t_0$
(frozen \emph{before} looking at any result, to avoid tuning it). For each
$t = t_0,\dots,T$:
\begin{enumerate}[nosep]
  \item form the estimation window $W_t = \{t-w,\dots,t-1\}$ (fixed length $w$,
        or expanding);
  \item fit the model on $W_t$ and produce $\widehat F_t$;
  \item record $\VaR_{\alpha,t}$, $\ES_{\alpha,t}$, $\sigma^2_{t+1}$, the
        realized $r_{t+1}$, the violation $I_{t+1}$ and the PIT $u_{t+1}$.
\end{enumerate}
No information dated after $t$ (and no data vintage newer than $t$) enters step 2.
Expensive models re-estimate parameters only every $k$ days while re-filtering
the state daily.

\intuition{The whole study is an honest simulation of running each model in
production: at every date it may only use the past, and the date at which we
stop cheating (train/test split) is chosen once and never moved.}

%======================================================================
\section{Realized measures from intraday data}\label{sec:realized}

Within day $t$ let $r_{t,1},\dots,r_{t,M}$ be the $M$ intraday
(5-minute) log-returns. All sums below are over $i$ within the day.

\paragraph{Realized variance and its subsampled version.}
\begin{equation}
  \RV_t = \sum_{i=1}^{M} r_{t,i}^2 .
\end{equation}
To damp microstructure noise we also report a subsampled estimator: aggregate
the 5-minute returns into $q$-step blocks starting at each of $q$ offsets and
average the resulting $\RV$'s.

\paragraph{Bipower variation and the jump component.}
With $\mu_1 = \E|Z| = \sqrt{2/\pi}$ for $Z\sim N(0,1)$,
\begin{equation}
  \BV_t = \mu_1^{-2}\,\frac{M}{M-1}\sum_{i=2}^{M} |r_{t,i}|\,|r_{t,i-1}|
        = \frac{\pi}{2}\,\frac{M}{M-1}\sum_{i=2}^{M} |r_{t,i}|\,|r_{t,i-1}| .
\end{equation}
$\BV_t$ is robust to jumps (a single large $|r_{t,i}|$ multiplies a
neighbouring \emph{small} return), so $\RV_t - \BV_t \ge 0$ isolates the jump
part of the quadratic variation. We gate it with the
Barndorff-Nielsen--Shephard (2006) $Z$-statistic
\begin{equation}
  Z_t \;=\; \frac{(\RV_t - \BV_t)/\RV_t}
       {\sqrt{\left(\tfrac{\pi^2}{4} + \pi - 5\right)\dfrac{1}{M}\,
        \max\!\left(1,\ \TQ_t/\BV_t^2\right)}},
  \qquad
  \text{jump}_t = (\RV_t - \BV_t)\,\Ind{Z_t > \Phi^{-1}(0.95)},
\end{equation}
where the tripower quarticity is
\begin{equation}
  \TQ_t = M\,\mu_{4/3}^{-3}\sum_{i=3}^{M} |r_{t,i}|^{4/3}|r_{t,i-1}|^{4/3}|r_{t,i-2}|^{4/3},
  \qquad
  \mu_{4/3} = 2^{2/3}\,\frac{\Gamma(7/6)}{\Gamma(1/2)}.
\end{equation}
An alternative per-bar detector (Lee--Mykland, 2008) flags $r_{t,i}$ as a jump
when $|r_{t,i}|/\hat\sigma_{t,i}$ exceeds a Gumbel critical value, with
$\hat\sigma_{t,i}$ a local bipower estimate.

\paragraph{Realized semivariance.}
\begin{equation}
  \RV^{+}_t = \sum_{i} r_{t,i}^2\,\Ind{r_{t,i} > 0},
  \qquad
  \RV^{-}_t = \sum_{i} r_{t,i}^2\,\Ind{r_{t,i} < 0},
  \qquad
  \RV^{+}_t + \RV^{-}_t = \RV_t .
\end{equation}
$\RV^{-}$ carries most of the predictive content for downside risk.

\paragraph{Realized quarticity.}
\begin{equation}
  \RQ_t = \frac{M}{3}\sum_{i=1}^{M} r_{t,i}^4 .
\end{equation}
$\RQ_t$ estimates the integrated quarticity $\int \sigma_s^4\,ds$; it is the
measurement-error variance of $\RV_t$ and is used by HARQ to down-weight noisy
days.

\intuition{Intraday data lets us \emph{observe} yesterday's variance instead of
inferring it from one squared return. $\RV$ is the observation, $\BV$ strips the
jumps, $\RV^{-}$ keeps the bad half, and $\RQ$ tells us how noisy the
observation itself is.}

%======================================================================
\section{The forecasting models}\label{sec:models}

Table~\ref{tab:models} lists the sixteen models by family. The rest of the
section gives each one's equations.

\begin{table}[t]
\centering
\small
\begin{tabular}{@{}lll@{}}
\toprule
Family & Models & Key idea \\
\midrule
Non-parametric        & HS, AWHS                 & empirical quantile of past returns \\
Exponential smoothing & EWMA                     & one-parameter recursive variance \\
Single-regime GARCH   & GARCH-$t$, GJR-$t$, EGARCH-$t$ & clustering $+$ leverage, $t$ innovations \\
Semiparametric tail   & FHS, GARCH-EVT           & GARCH filter $+$ empirical / GPD residual tail \\
Discontinuous         & Merton jump-diffusion    & Brownian part $+$ compound-Poisson jumps \\
Realized-measure      & HAR-RV, HARQ, Realized GARCH & forecast from observed $\RV$ \\
Exogenous / quantile  & GARCH-X, CAViaR, CAViaR-X & driver in the variance / quantile eqn.\\
Regime-switching      & MS-GARCH                  & 2-state Markov mixture of GARCH \\
\bottomrule
\end{tabular}
\caption{The model roster.}
\label{tab:models}
\end{table}

%----------------------------------------------------------------------
\subsection{Historical Simulation and age-weighting}

Plain HS takes the empirical $\alpha$-quantile and tail mean of the window
$\{r_s\}_{s\in W_t}$:
\begin{equation}
  \VaR_{\alpha,t} = \widehat Q_\alpha\big(\{r_s\}_{s\in W_t}\big),
  \qquad
  \ES_{\alpha,t} = \frac{\sum_{s\in W_t} r_s\,\Ind{r_s \le \VaR_{\alpha,t}}}
                        {\sum_{s\in W_t} \Ind{r_s \le \VaR_{\alpha,t}}} .
\end{equation}
Age-weighted HS (Boudoukh--Richardson--Whitelaw, 1998) replaces the uniform
weights by geometrically decaying ones. With half-life $H$ and
$\lambda = 2^{-1/H}$, observation $s$ at age $a_s = t-1-s$ gets
$w_s \propto \lambda^{\,a_s}$, $\sum_s w_s = 1$, and $\widehat Q_\alpha$ is the
weighted quantile: the smallest $r_{(k)}$ with $\sum_{j\le k} w_{(j)} \ge \alpha$
(order statistics).

\intuition{``What has actually happened lately, in order of severity.'' No
distributional assumption; age-weighting just says last month matters more than
last year.}

%----------------------------------------------------------------------
\subsection{EWMA / RiskMetrics}

Zero conditional mean and a one-parameter variance recursion:
\begin{equation}
  \sigma^2_{t+1} = \lambda\,\sigma^2_{t} + (1-\lambda)\,r_t^2,
  \qquad \lambda = 0.94 .
\end{equation}
The predictive law is $r_{t+1} \sim \sigma_{t+1} Z$ with $Z\sim N(0,1)$ (or a
fixed-$\nu$ standardized $t$), so
$\VaR_{\alpha,t} = \sigma_{t+1}\,\Phi^{-1}(\alpha)$ and
$\ES_{\alpha,t} = -\sigma_{t+1}\,\phi(\Phi^{-1}(\alpha))/\alpha$.

\intuition{An exponentially fading average of squared returns. Cheap, hard to
beat on level, but the Gaussian tail is too thin for crypto.}

%----------------------------------------------------------------------
\subsection{Single-regime GARCH with standardized \texorpdfstring{$t$}{t} innovations}

Mean $\mu$ constant; $r_t = \mu + \sqrt{h_t}\,\eta_t$, $\eta_t \sim
\text{standardized } t_\nu$ (unit variance). The three variants differ only in
the variance recursion:
\begin{align}
  \text{GARCH-}t:\quad & h_t = \omega + \alpha\,\varepsilon_{t-1}^2 + \beta\,h_{t-1}, \\
  \text{GJR-}t:\quad   & h_t = \omega + \big(\alpha + \gamma\,\Ind{\varepsilon_{t-1}<0}\big)\varepsilon_{t-1}^2 + \beta\,h_{t-1}, \\
  \text{EGARCH-}t:\quad& \ln h_t = \omega + \alpha\big(|\eta_{t-1}| - \E|\eta_{t-1}|\big) + \gamma\,\eta_{t-1} + \beta\,\ln h_{t-1},
\end{align}
with $\varepsilon_t = r_t - \mu$. Parameters are estimated by ML on the window
(returns scaled $\times 100$ for the optimiser).

\paragraph{Standardization of the $t$ quantile (this is the bug v1 got wrong).}
A raw Student-$t_\nu$ variable has variance $\nu/(\nu-2)$. The fitted $\eta_t$
has \emph{unit} variance, so its $\alpha$-quantile is
\begin{equation}
  z_\alpha = \sqrt{\tfrac{\nu-2}{\nu}}\; t_\nu^{-1}(\alpha),
\end{equation}
and, writing $q_\alpha = t_\nu^{-1}(\alpha)$ and $f_\nu$ for the raw-$t$ density,
its lower-tail mean is
\begin{equation}
  \bar z_\alpha
    = -\sqrt{\tfrac{\nu-2}{\nu}}\;\frac{f_\nu(q_\alpha)}{\alpha}\,
      \frac{\nu + q_\alpha^2}{\nu - 1}.
\end{equation}
The one-step forecasts are
$\VaR_{\alpha,t} = \mu + \sqrt{h_{t+1}}\,z_\alpha$ and
$\ES_{\alpha,t} = \mu + \sqrt{h_{t+1}}\,\bar z_\alpha$.
Omitting $\sqrt{(\nu-2)/\nu}$ inflates VaR by $\sqrt{\nu/(\nu-2)}$ ($\approx
1.7\times$ at $\nu = 3$). EGARCH's log-variance forecast can also blow up near
the non-stationary boundary; when $\sqrt{h_{t+1}}$ exceeds $20\times$ the window
standard deviation the model falls back to the empirical quantile.

\intuition{GARCH says today's variance is a weighted memory of yesterday's shock
and yesterday's variance; GJR/EGARCH add that a bad shock hurts more than a good
one; the $t$ replaces the Gaussian tail with a fat one whose thickness $\nu$ is
estimated.}

%----------------------------------------------------------------------
\subsection{Filtered Historical Simulation (FHS)}

Barone-Adesi--Giannopoulos--Vosper (1999): fit a GARCH(1,1) filter, take the
standardized residuals $\hat z_s = \hat\varepsilon_s / \sqrt{\hat h_s}$,
$s \in W_t$, and rescale their empirical quantile by the one-step volatility:
\begin{equation}
  \VaR_{\alpha,t} = \hat\mu + \sqrt{\hat h_{t+1}}\;\widehat Q_\alpha(\{\hat z_s\}),
  \qquad
  \ES_{\alpha,t}  = \hat\mu + \sqrt{\hat h_{t+1}}\;
       \frac{\sum_s \hat z_s\,\Ind{\hat z_s \le \widehat Q_\alpha}}
            {\sum_s \Ind{\hat z_s \le \widehat Q_\alpha}} .
\end{equation}

\intuition{Let GARCH remove the volatility dynamics, then read the tail off the
leftover shocks without pretending they are Gaussian. The residual shape is
whatever the data says.}

%----------------------------------------------------------------------
\subsection{GARCH-EVT (McNeil--Frey, 2000)}

Same GARCH filter as FHS, but the lower tail of the standardized residuals is
modelled with a Generalized Pareto Distribution by Peaks-Over-Threshold. Work in
loss units $L = -z$. Fix a threshold $u$ (the $90\%$ empirical loss quantile),
let $N_u$ of the $n$ residuals exceed $u$, and fit
$(\xi,\beta)$ to the exceedances $\{L_s - u : L_s > u\}$ by ML. The tail
approximation
\begin{equation}
  \Prob(L > x) \approx \frac{N_u}{n}\left(1 + \xi\,\frac{x-u}{\beta}\right)^{-1/\xi},
  \qquad x > u,
\end{equation}
inverts, for $\alpha < N_u/n$, to a residual VaR and ES in loss units
\begin{equation}
  \VaR^{L}_\alpha = u + \frac{\beta}{\xi}\!\left[\left(\frac{\alpha}{N_u/n}\right)^{-\xi} - 1\right],
  \qquad
  \ES^{L}_\alpha = \frac{\VaR^{L}_\alpha}{1-\xi} + \frac{\beta - \xi u}{1-\xi}
                   \quad (\xi < 1),
\end{equation}
(with the $\xi\to 0$ exponential limit $\VaR^{L}_\alpha = u + \beta\ln\frac{N_u/n}{\alpha}$).
Back in return units, $z$-quantile $= -\VaR^{L}_\alpha$, and
$\VaR_{\alpha,t} = \hat\mu + \sqrt{\hat h_{t+1}}\,(-\VaR^{L}_\alpha)$, similarly
for ES. When $\alpha \ge N_u/n$ (the quantile is inside the body) or the fit
degenerates, the model falls back to the empirical residual quantile. The body
of the residual CDF (for PIT) is the empirical CDF; the tail is the GPD
expression above.

\intuition{Extreme value theory gives the only asymptotically justified shape
for ``losses beyond a high threshold.'' We fit that shape to the GARCH-filtered
shocks, so the tail is estimated from the tail, not extrapolated from the
centre.}

%----------------------------------------------------------------------
\subsection{Merton jump-diffusion}

Daily returns are a Brownian part plus a compound-Poisson jump part
(Merton, 1976):
\begin{equation}
  r_{t+1} = \Big(\mu - \tfrac12\sigma^2 - \lambda\bar k\Big)
            + \sigma\,\varepsilon_{t+1}
            + \sum_{j=1}^{N_{t+1}} J_j,
  \qquad
  N_{t+1}\sim \text{Poisson}(\lambda),\quad J_j \sim N(\mu_J,\sigma_J^2),
\end{equation}
with drift compensator $\bar k = e^{\mu_J + \sigma_J^2/2} - 1$. The density of
$r_{t+1}$ is a Poisson mixture of normals,
\begin{equation}
  f(r) = \sum_{n=0}^{\infty}
         \frac{e^{-\lambda}\lambda^n}{n!}\,
         \varphi\!\big(r;\ \mu - \lambda\bar k + n\mu_J,\ \sigma^2 + n\sigma_J^2\big),
\end{equation}
truncated at $n = 10$ for the log-likelihood, which is maximised over
$(\lambda,\mu_J,\sigma_J)$ on the window (the sum is vectorised over $n$; $\mu$
and $\sigma$ are the sample moments). The predictive law is represented by a
large Monte-Carlo sample of $r_{t+1}$, from which $\VaR$, $\ES$ and the CDF are
read; $\sigma^2_{t+1} = \sigma^2 + \lambda(\mu_J^2 + \sigma_J^2)$.

\intuition{Some crypto down-days are not ``a big draw from a fat distribution''
but a discrete event (an exchange failure, a de-peg). The Poisson term models
those directly, on a different time scale from the GARCH clustering.}

%----------------------------------------------------------------------
\subsection{HAR-RV and HARQ}

Corsi (2009): regress the realized variance on its own daily, weekly and
monthly averages,
\begin{equation}
  \RV_{t} = \beta_0
          + \beta_d\,\RV_{t-1}
          + \beta_w\,\overline{\RV}^{(w)}_{t-1}
          + \beta_m\,\overline{\RV}^{(m)}_{t-1}
          + \varepsilon_{t},
  \qquad
  \overline{\RV}^{(w)}_{t-1} = \tfrac{1}{5}\!\sum_{i=1}^{5}\RV_{t-i},\quad
  \overline{\RV}^{(m)}_{t-1} = \tfrac{1}{22}\!\sum_{i=1}^{22}\RV_{t-i},
\end{equation}
estimated by OLS on the window. HARQ (Bollerslev--Patton--Quaedvlieg, 2016)
interacts the daily lag with realized quarticity, so the model trusts the daily
term less on noisy days:
\begin{equation}
  \RV_{t} = \beta_0
          + \Big(\beta_d + \beta_{dQ}\,\widetilde{\RQ}^{\,1/2}_{t-1}\Big)\RV_{t-1}
          + \beta_w\,\overline{\RV}^{(w)}_{t-1}
          + \beta_m\,\overline{\RV}^{(m)}_{t-1}
          + \varepsilon_{t},
\end{equation}
where $\widetilde{\RQ}^{\,1/2}$ is $\RQ^{1/2}$ standardized (subtract mean,
divide by s.d.) so the interaction sits on the same scale as the other
regressors. The one-step forecast $\widehat\RV_{t+1}$ \emph{is}
$\sigma^2_{t+1}$; it is floored at half the $5\%$ RV quantile and capped at
$30\times$ the median RV (the $\RQ$ interaction can otherwise explode on an
extreme-quarticity day). The return law is $\sqrt{\sigma^2_{t+1}}$ times a
standardized $t_\nu$, with $\nu$ fitted to $r_s/\sqrt{\RV_s}$ on the window.

\intuition{Volatility is persistent at several horizons at once --- today, this
week, this month --- and a linear model in those three averages captures most of
it. HARQ adds: on a day when the variance estimate itself was noisy, lean less
on it.}

%----------------------------------------------------------------------
\subsection{Realized GARCH (Hansen--Huang--Shek, 2012)}

A GARCH whose latent variance is tied to the observed $\RV$ through a
measurement equation. With $z_t = r_t/\sqrt{h_t}$, $z_t \sim N(0,1)$,
$x_t = \RV_t$:
\begin{align}
  r_t &= \sqrt{h_t}\,z_t, \\
  \ln h_t &= \omega + \beta\,\ln h_{t-1} + \gamma\,\ln x_{t-1}, \\
  \ln x_t &= \xi + \varphi\,\ln h_t + \tau_1 z_t + \tau_2(z_t^2 - 1) + u_t,
            \qquad u_t \sim N(0,\sigma_u^2).
\end{align}
The leverage function $\tau(z) = \tau_1 z + \tau_2(z^2-1)$ lets the measurement
depend asymmetrically on the return sign. The eight parameters are estimated by
joint Gaussian ML,
\begin{equation}
  \ell = -\tfrac12\sum_t\big(\ln h_t + z_t^2\big)
         -\tfrac12\sum_t\Big(\ln(2\pi\sigma_u^2) + u_t^2/\sigma_u^2\Big),
\end{equation}
and the $\ln h_t$ recursion, being AR(1) with a time-varying input, is solved in
closed form (a linear filter) rather than a Python loop. Forecast:
$h_{t+1} = \exp\!\big(\omega + \beta\ln h_t + \gamma\ln x_t\big)$.

\intuition{Ordinary GARCH updates variance from the squared return, a very noisy
signal. Realized GARCH updates it from $\RV$ instead, and the measurement
equation keeps the two definitions of ``variance'' mutually consistent.}

%----------------------------------------------------------------------
\subsection{GARCH-X}

GARCH(1,1)-$t$ with a lagged realized variance in the variance equation:
\begin{equation}
  h_t = \omega + \alpha\,r_{t-1}^2 + \beta\,h_{t-1} + \gamma\,x_{t-1},
  \qquad r_t/\sqrt{h_t} \sim \text{standardized } t_\nu,\quad x_t = \RV_t .
\end{equation}
Estimated by ML (\texttt{arch} has no variance-equation exogenous term, so this
is hand-coded); the recursion is again a linear filter. Without a realized block
it reduces to plain GARCH-$t$ ($\gamma = 0$). The same $20\times$ volatility
ceiling as the GARCH family applies.

\intuition{A minimal way to ask ``does yesterday's \emph{measured} variance add
anything on top of yesterday's squared return?'' --- the coefficient $\gamma$
answers it.}

%----------------------------------------------------------------------
\subsection{CAViaR and CAViaR-X (Engle--Manganelli, 2004)}

Model the conditional quantile directly, with no distribution. Two
specifications:
\begin{align}
  \text{SAV:}\quad & q_t = \beta_1 + \beta_2\,q_{t-1} + \beta_3\,|r_{t-1}|, \\
  \text{AS:}\quad  & q_t = \beta_1 + \beta_2\,q_{t-1}
                    + \beta_3\,(r_{t-1})^{+} + \beta_4\,(r_{t-1})^{-},
\end{align}
with $(x)^{+}=\max(x,0)$, $(x)^{-}=\max(-x,0)$; CAViaR-X adds $\beta_5\,x_{t-1}$
with $x_{t-1} = \sqrt{\RV_{t-1}}$. Given the window, parameters minimise the
\emph{tick} (pinball) loss at level $\alpha$,
\begin{equation}
  \hat\theta = \arg\min_\theta \ \frac{1}{n}\sum_{s\in W_t}
     \big(\alpha - \Ind{r_s < q_s(\theta)}\big)\big(r_s - q_s(\theta)\big),
\end{equation}
whose population minimiser is exactly the $\alpha$-quantile. The recursion
$q_t = \beta_1 + \beta_2 q_{t-1} + (\text{news})_{t-1}$ is AR(1) with a
time-varying input and is evaluated with a linear filter, making the
optimisation fast. CAViaR is fit \emph{once per $\alpha$}; a crossing repair
enforces $q^{(\alpha')} \le q^{(\alpha)}$ for $\alpha' < \alpha$. Expected
Shortfall has no direct CAViaR analogue; we scale the forecast quantile by the
in-sample tail-to-quantile ratio,
\begin{equation}
  \widehat{\ES}_{\alpha,t} = q_{t+1}\cdot\hat\rho,
  \qquad
  \hat\rho = \frac{\overline{r_s \mid r_s < q_s}}{\overline{q_s \mid r_s < q_s}}
             \in [1,\,3]\ \text{(clipped)} .
\end{equation}

\intuition{Skip the density entirely: write down how the $2.5\%$ worst outcome
evolves --- it is persistent (autoregressive) and it widens after big moves ---
and fit that curve to make the empirical breach rate come out at $2.5\%$.}

%----------------------------------------------------------------------
\subsection{MS-GARCH (2-state Markov-switching)}

A latent state $S_t \in \{1,2\}$ follows a Markov chain with transition matrix
$P = (p_{k\ell})$. Conditional on the state, the variance is a separate
GARCH(1,1) with $t$ innovations:
\begin{equation}
  r_t \mid (\Fcal_{t-1}, S_t = k) \sim \sqrt{h_{k,t}}\;\eta_{k,t},
  \qquad
  h_{k,t} = \omega_k + \alpha_k\,r_{t-1}^2 + \beta_k\,h_{k,t-1},
  \qquad
  \eta_{k,t}\sim \text{standardized } t_{\nu_k}.
\end{equation}
Let $\pi_{k,t} = \Prob(S_t = k \mid \Fcal_{t-1})$ be the one-step predictive
regime probabilities (Hamilton filter). The predictive density of $r_{t+1}$ is
the mixture
\begin{equation}
  f_{t}(r) = \sum_{k=1}^{2} \pi_{k,t+1}\; f_k\!\big(r;\ h_{k,t+1},\ \nu_k\big),
\end{equation}
and $\VaR_{\alpha,t}$, $\ES_{\alpha,t}$ are obtained numerically from its CDF.
Estimation is by ML in R (\texttt{MSGARCH}); the walk-forward re-fits every $20$
days and re-filters the state daily through appended observations. Regimes are
labelled by stationary variance
$\omega_k/(1-\alpha_k-\beta_k)$ each re-fit to avoid label switching.
The bridge model in Python serves the cached R predictions; dates without a
cached value fall back to the empirical quantile.

The filtered probability $\Prob(S_t = k^\star \mid \Fcal_t)$ of the high-variance
regime is also recorded. A single full-sample fit's filtered probability
correlates $\approx 0.7$ with $|r_t|$; the rolling-window walk-forward version
does not ($\approx 0.07$), because a two-regime GARCH is weakly identified on
$\sim 500$ observations. The descriptive regime curve therefore uses the
full-sample fit and is kept out of the VaR backtest.

\intuition{Markets alternate between a calm regime and a turbulent one, and
which one you are in is not observed. The model carries a probability over both
and mixes their forecasts. It works as a density; extracting a real-time
``are we in a crisis'' signal from short windows does not.}

%======================================================================
\section{Backtesting and comparative evaluation}\label{sec:eval}

Two questions. First, \emph{is each model correctly specified?} --- coverage,
independence, the Dynamic Quantile test, and Expected-Shortfall backtests.
Second, \emph{which models forecast best, and is the difference real?} ---
strictly consistent scoring functions, the Diebold--Mariano test, and the Model
Confidence Set. Berkowitz closes the loop on the whole density.

Throughout, $n$ is the number of out-of-sample days, $x = \sum_t I_t$ the number
of violations, $\hat\pi = x/n$, and the rejection level for the specification
tests is $5\%$.

%----------------------------------------------------------------------
\subsection{Unconditional coverage --- Kupiec POF (1995)}

$H_0$: the violation rate equals $\alpha$. The likelihood-ratio statistic
\begin{equation}
  \mathrm{LR}_{uc}
   = -2\ln\frac{(1-\alpha)^{\,n-x}\,\alpha^{\,x}}
               {(1-\hat\pi)^{\,n-x}\,\hat\pi^{\,x}}
   \ \xrightarrow{d}\ \chi^2_1 .
\end{equation}

\intuition{Did the $2.5\%$ line get crossed about $2.5\%$ of the time? Nothing
about \emph{when}.}

%----------------------------------------------------------------------
\subsection{Independence and conditional coverage --- Christoffersen (1998)}

Let $n_{ij}$ count transitions $I_{t-1} = i \to I_t = j$, with
$\hat\pi_{ij} = n_{ij}/(n_{i0}+n_{i1})$ and
$\hat\pi = (n_{01}+n_{11})/n$. The independence LR against a first-order Markov
alternative is
\begin{equation}
  \mathrm{LR}_{ind}
   = -2\ln\frac{(1-\hat\pi)^{\,n_{00}+n_{10}}\;\hat\pi^{\,n_{01}+n_{11}}}
               {(1-\hat\pi_{01})^{\,n_{00}}\hat\pi_{01}^{\,n_{01}}\,
                (1-\hat\pi_{11})^{\,n_{10}}\hat\pi_{11}^{\,n_{11}}}
   \ \xrightarrow{d}\ \chi^2_1 ,
\end{equation}
and conditional coverage combines the two:
$\mathrm{LR}_{cc} = \mathrm{LR}_{uc} + \mathrm{LR}_{ind} \xrightarrow{d} \chi^2_2$.

\intuition{If a violation today makes a violation tomorrow more likely, the
model is missing volatility clustering even if the average rate is fine.}

%----------------------------------------------------------------------
\subsection{Dynamic Quantile test --- Engle--Manganelli (2004)}

The most discriminating coverage test. Form the demeaned hit
$\mathrm{Hit}_t = I_t - \alpha$ and regress it on a constant, its own $p$ lags
and the contemporaneous VaR:
\begin{equation}
  \mathrm{Hit}_t = \delta_0
    + \sum_{j=1}^{p}\delta_j\,\mathrm{Hit}_{t-j}
    + \delta_{p+1}\,\VaR_{\alpha,t}
    + \epsilon_t .
\end{equation}
Under correct specification $\E[\mathrm{Hit}_t \mid \Fcal_{t-1}] = 0$, so every
$\delta$ is zero and, with $X$ the regressor matrix,
\begin{equation}
  \mathrm{DQ} = \frac{\hat\delta^{\top}(X^{\top}X)\,\hat\delta}{\alpha(1-\alpha)}
   \ \xrightarrow{d}\ \chi^2_{p+2}
   \qquad (\text{we use } p = 4).
\end{equation}

\intuition{Kupiec and Christoffersen only look at the hit series; DQ also asks
whether the hits can be \emph{predicted} from recent hits or from the VaR level
itself. If a model breaches mostly when its own VaR is shallow, DQ catches it.}

%----------------------------------------------------------------------
\subsection{Basel traffic light}

Count exceptions in a rolling $250$-day window: $\le 4$ green,
$5$--$9$ amber, $\ge 10$ red, mapping to a capital-multiplier add-on
$0,\ 0.40\text{--}0.85,\ 1.0$. A regulatory, not statistical, check; used in the
capital layer.

%----------------------------------------------------------------------
\subsection{Expected-Shortfall backtests --- Acerbi--Székely (2014)}

ES is not elicitable on its own, so these are diagnostics rather than scores.
With $v_t = \VaR_{\alpha,t} < 0$, $e_t = \ES_{\alpha,t} < 0$, breaches
$\mathcal{B} = \{t : r_t < v_t\}$:
\begin{equation}
  Z_2 = -\frac{1}{n\alpha}\sum_{t\in\mathcal B}\frac{r_t}{e_t} + 1,
  \qquad
  Z_1 = -\frac{1}{|\mathcal B|}\sum_{t\in\mathcal B}\frac{r_t}{e_t} + 1 .
\end{equation}
Under $H_0$ (ES correctly specified) $\E[Z_1] = \E[Z_2] = 0$; $Z_2$ also folds
in the breach frequency, $Z_1$ conditions on being in the tail. $Z < 0$ means
realized tail losses are worse than the ES forecast. The null has no clean
closed form; the $p$-value is
$\widehat p = \tfrac{1}{B}\sum_{b} \Ind{Z^{(b)} \le Z^{\text{obs}}}$ over $B$
return paths simulated from the model's own predictive distribution each day.

\intuition{Given the model says ``when you breach, you lose about $e_t$ on
average,'' compare that to how much you \emph{actually} lost on breach days. If
the average realized breach is $30\%$ deeper than $e_t$, $Z_2$ is clearly
negative.}

%----------------------------------------------------------------------
\subsection{Strictly consistent scoring: the FZ0 loss}

The Fissler--Ziegel family gives losses that are minimised in expectation
\emph{exactly} at the true $(\VaR,\ES)$ pair; the $0$-homogeneous member
(Patton--Ziegel--Chen, 2019), with $v_t, e_t < 0$, is
\begin{equation}
  L^{\mathrm{FZ0}}_t
   = \frac{\Ind{r_t \le v_t}}{\alpha\,e_t}\,(r_t - v_t)
   + \frac{v_t}{e_t} + \ln(-e_t) - 1 .
\end{equation}
The breach term is positive ($1/(\alpha e_t) < 0$ and $r_t - v_t < 0$), so a
too-shallow VaR with many breaches is penalised. Lower average $L^{\mathrm{FZ0}}$
is unambiguously better, which is what makes model \emph{ranking} legitimate.
For variance forecasts alone we use QLIKE (Patton, 2011),
\begin{equation}
  L^{\mathrm{QL}}_t = \frac{\RV_{t}}{\sigma^2_{t}} - \ln\frac{\RV_{t}}{\sigma^2_{t}} - 1
  \;\ge\; 0,
\end{equation}
which is robust to the fact that $\RV$ is a noisy proxy for the true variance.

\intuition{A consistent loss is a scoring rule you cannot game: the way to make
it small \emph{is} to report the true VaR and ES. Averaging it over the
out-of-sample period turns ``which model is better'' into a number.}

%----------------------------------------------------------------------
\subsection{Diebold--Mariano (1995)}

For two models with loss difference $d_t = L^A_t - L^B_t$, test
$H_0: \E[d_t] = 0$. With a Newey--West long-run variance
\begin{equation}
  \widehat{\mathrm{Var}}(\bar d)
   = \frac{1}{n}\left(\hat\gamma_0
     + 2\sum_{k=1}^{L}\Big(1 - \tfrac{k}{L+1}\Big)\hat\gamma_k\right),
  \qquad L \approx n^{1/3},
\end{equation}
the statistic $\mathrm{DM} = \bar d / \sqrt{\widehat{\mathrm{Var}}(\bar d)}$ is
multiplied by the Harvey--Leybourne--Newbold small-sample factor
$\sqrt{\big(n + 1 - 2L + L(L-1)/n\big)/n}$ and referred to $t_{n-1}$.

\intuition{One model's average loss being lower is not enough --- the
difference has to be large relative to how much it wobbles over time (and
loss differences are autocorrelated, hence the HAC variance).}

%----------------------------------------------------------------------
\subsection{Model Confidence Set --- Hansen--Lunde--Nason (2011)}

Rather than pick one winner, isolate the set of models that are statistically
indistinguishable from the best. Let $\mathcal M$ be the current set,
$L_{i,t}$ the loss of model $i$. Define the excess loss relative to the set
mean and its time average,
\begin{equation}
  d_{i,t} = L_{i,t} - \frac{1}{|\mathcal M|}\sum_{j\in\mathcal M} L_{j,t},
  \qquad
  \bar d_i = \frac{1}{T}\sum_{t} d_{i,t},
  \qquad
  t_i = \frac{\bar d_i}{\sqrt{\widehat{\mathrm{var}}^{*}(\bar d_i)}},
\end{equation}
with $\widehat{\mathrm{var}}^{*}$ from a stationary block bootstrap (block
lengths $\sim \mathrm{Geom}(1/\ell)$, wrap-around resampling). The equivalence
statistic is $T_{\max} = \max_{i} t_i$, and its null distribution is the
bootstrap distribution of
$\max_i (\,\bar d_i^{*} - \bar d_i\,)/\sqrt{\widehat{\mathrm{var}}^{*}(\bar d_i)}$.
The algorithm:
\begin{enumerate}[nosep]
  \item compute $T_{\max}$ and its bootstrap $p$-value $p$;
  \item set $p_{\text{run}} \leftarrow \max(p_{\text{run}}, p)$; if
        $p_{\text{run}} \ge 1 - \text{confidence}$, stop --- every model left is
        in the MCS;
  \item otherwise remove $\arg\max_i t_i$ (the worst), record its MCS $p$-value
        as $p_{\text{run}}$, and repeat.
\end{enumerate}
The output is the surviving set (ordered by average loss) and, for every model,
its MCS $p$-value. This is the headline result of the study: it is run on the
FZ0 losses (the $(\VaR,\ES)$ ranking) and separately on the QLIKE losses (the
variance ranking), per asset and per $\alpha$.

\intuition{Like a tournament with confidence intervals. Repeatedly ask ``do the
remaining models all forecast equally well?''; while the answer is no, drop the
worst; when the answer becomes yes, whoever is left is the confidence set. With
little data the set is large; with a lot it narrows to the genuine leaders.}

%----------------------------------------------------------------------
\subsection{Density calibration --- Berkowitz (2001)}

If the whole predictive density is right, the PIT $u_t = \widehat F_{t-1}(r_t)$
is i.i.d.\ $\mathrm{U}(0,1)$, hence $z_t = \Phi^{-1}(u_t)$ is i.i.d.\ $N(0,1)$.
Fit an AR(1) $z_t = \mu + \rho\,z_{t-1} + \epsilon_t$,
$\epsilon_t \sim N(0,\sigma^2)$, and LR-test
$H_0: (\mu,\rho,\sigma^2) = (0,0,1)$:
\begin{equation}
  \mathrm{LR}_{B} = -2\big(\ell(0,0,1) - \ell(\hat\mu,\hat\rho,\hat\sigma^2)\big)
   \ \xrightarrow{d}\ \chi^2_3 .
\end{equation}

\intuition{One test for the whole forecast distribution, not just the tail:
after mapping outcomes through the model's own CDF, what is left should look
like independent standard-normal noise --- no drift, no memory, right spread.}

%======================================================================
\section*{A note on what the mathematics does not fix}

The tools above are honest about three things the crypto data will not give up.
(i) A two-regime MS-GARCH is not identified on rolling $500$-day windows, so its
real-time regime probability carries no out-of-sample signal --- only the
full-sample fit does, and that is descriptive. (ii) HAR-type models forecast the
\emph{mean} of $\RV$, which understates the volatility relevant on a tail day,
so they tend to run a little thin for VaR. (iii) EWMA and EGARCH keep tails too
thin for $\alpha = 0.01$. None of these are bugs; they are properties, and the
FZ0/MCS machinery is precisely what surfaces them without hand-waving.

%======================================================================
\section*{References (abbreviated)}

\begin{itemize}[nosep,leftmargin=*]
\item Acerbi \& Székely (2014). Back-testing Expected Shortfall. \emph{Risk}.
\item Barndorff-Nielsen \& Shephard (2006). Econometrics of testing for jumps. \emph{J.\ Fin.\ Econometrics}.
\item Barone-Adesi, Giannopoulos \& Vosper (1999). VaR without correlations for portfolios of derivative securities. \emph{J.\ Futures Markets}.
\item Berkowitz (2001). Testing density forecasts. \emph{JBES}.
\item Bollerslev, Patton \& Quaedvlieg (2016). Exploiting the errors: a simple approach for improved volatility forecasting. \emph{J.\ Econometrics}.
\item Boudoukh, Richardson \& Whitelaw (1998). The best of both worlds. \emph{Risk}.
\item Christoffersen (1998). Evaluating interval forecasts. \emph{Int.\ Econ.\ Review}.
\item Corsi (2009). A simple approximate long-memory model of realized volatility. \emph{J.\ Fin.\ Econometrics}.
\item Diebold \& Mariano (1995). Comparing predictive accuracy. \emph{JBES}.
\item Engle \& Manganelli (2004). CAViaR. \emph{JBES}.
\item Hansen, Huang \& Shek (2012). Realized GARCH. \emph{J.\ Applied Econometrics}.
\item Hansen, Lunde \& Nason (2011). The Model Confidence Set. \emph{Econometrica}.
\item Harvey, Leybourne \& Newbold (1997). Testing the equality of prediction mean squared errors. \emph{Int.\ J.\ Forecasting}.
\item Kupiec (1995). Techniques for verifying the accuracy of risk measurement models. \emph{J.\ Derivatives}.
\item Lee \& Mykland (2008). Jumps in financial markets. \emph{Review of Financial Studies}.
\item McNeil \& Frey (2000). Estimation of tail-related risk measures for heteroscedastic financial time series. \emph{J.\ Empirical Finance}.
\item Merton (1976). Option pricing when underlying stock returns are discontinuous. \emph{J.\ Fin.\ Economics}.
\item Patton (2011). Volatility forecast comparison using imperfect volatility proxies. \emph{J.\ Econometrics}.
\item Patton, Ziegel \& Chen (2019). Dynamic semiparametric models for expected shortfall (and value-at-risk). \emph{J.\ Econometrics}.
\end{itemize}

\end{document}
